\section{Why a two-sided point is of type I}

Write $F$ for the original marked arithmetic frame, $O$ for its
coefficient DVR, $\eta$ for its sign and $\Psi$ for its character.
Write $c_*$ for its distinguished involution and $k$ for its residue
field. Indices of the three matrix rows and columns are $0,1,2$.
The block alternative of the shape calculation would leave a
two-dimensional irreducible representation on coordinates $0,2$.
Its lower-left entry is detected by the weight $-2$ Selmer group.

% BEGIN TSI1
\paragraph{TSI1. A normalized residual cocycle in weight minus two.}
Let $p(g)=\eta(g)\bar\Psi(g)^{-2}$. Let $x$ be a locally constant
$k$-valued cocycle for $p$, zero at $c_*$. Suppose it vanishes on
the specified $N_v$ away from $T_b$, and at every $v\in T_b$ is
of the form $(p(g)-1)w_v$ on the kernel of the local sign.
Under the frame conditions and $F.\mathrm{Xpp}$, $x=0$.

Insert $x$ in entry $(2,0)$ of the weight $-2$ matrix space;
the group action there is exactly multiplication by $p(g)$.
On $N_v$ in $T_b$, both the sign and the residual character are
trivial, so the cocycle vanishes there as well. At $v_\ell$ the
local Selmer condition asks for a coboundary on the whole group.
Subtract the displayed coboundary. The difference $a$ vanishes on
the local sign kernel. Within that kernel the frame provides $s$
with $\bar\Psi(s)^2\ne1$. For any local $d$ compute $a(ds)$ first
as $a(d)+p(d)a(s)$ and then as
$a(dsd^{-1})+p(dsd^{-1})a(d)$. This gives
$(1-p(s))a(d)=0$, so the difference vanishes. The weight-zero
extra condition is inapplicable in degree $-2$. The matrix cocycle
now lies in the specified Selmer space. Its quotient by coboundaries
vanishes by $F.\mathrm{Xpp}$, hence $x(g)=(p(g)-1)m$ globally.
At $c_*$ we have $p(c_*)=-1$ and $x(c_*)=0$. Since $2$ is a unit,
$m=0$. The same conclusion holds over any residue field extension
if the local coboundaries and values are taken there: apply each
$k$-linear functional, then use separation by the linear dual.
% END TSI1

% BEGIN TSI2
\paragraph{TSI2. A lower-left entry cannot survive the block condition.}
Let $M_g$ be the integral matrix of a marked point over a finite
coefficient DVR $D$ with residue $k'$. Suppose $(M_g)_{21}=0$
on $\ker\eta$. It vanishes on all of $\Gamma$: outside the kernel
use $g=(gc_*)c_*$ and $M_{c_*}=1$. If some $(M_g)_{20}$ were nonzero,
choose $g_0$ with minimal valuation of its nonzero entry, put
$a=(M_{g_0})_{20}$ and write $(M_g)_{20}=a b_g$. Then $b_{g_0}=1$.
Continuity modulo $\mathfrak m_D^{v(a)+1}$ of the original entry
gives local constancy of $\bar b_g$. Define $z_g=\bar b_g\bar\Psi(g)^{-1}$.

For $\eta(g)=1$, the $(2,0)$ entry of $M_gM_h$ after division by
$a$ gives $b_g(M_h)_{00}+(M_g)_{22}b_h$. For negative sign the
multiplication law uses the reversed transpose inverse of $M_h$.
The two inverse entries required here have numerators
$(M_h)_{00}(M_h)_{11}-(M_h)_{01}(M_h)_{10}$ and
$-(M_h)_{11}a b_h$, respectively, since $(M_h)_{21}=0$.
The determinant reduces to $1$. Using the diagonal residual entries
$\bar\Psi,1,\bar\Psi^{-1}$, both signs give
\[
 \bar b_{gh}=\bar b_g\bar\Psi(h)
             +\eta(g)\bar\Psi(g)^{-1}\bar b_h,\qquad
 z_{gh}=z_g+p(g)z_h .
\]
The normalized value at $c_*$ is zero. Away from $T_b$ the original
point is trivial on $N_v$, giving zero there.

For $v\in T_b$ take a local upper-triangularizing basis
$H\equiv 1\bmod\mathfrak m_D$ provided by the ordinary or diagonal
condition. For positive-sign $g$ the first column of $H$ spans an
invariant line. Eliminate its eigenvalue between rows $0$ and $2$
of $M_gH e_0$ to get
\[
 a b_g H_{00}^2=H_{20} U_g,\quad
 U_g=(M_gH)_{00}-H_{00}(M_g)_{22},\qquad
 \bar U_g=\bar\Psi(g)-\bar\Psi(g)^{-1}.
\]
Choose the frame separator $s$ there, so $\bar U_s\ne0$.
Cross-multiply the equations for $g,s$, cancel $aH_{00}^2$,
and reduce. It follows that $\bar b_g=\bar U_g\,w$ with
$w=\bar b_s/\bar U_s$. Consequently $z_g=(p(g)-1)(-w)$ on the
local sign kernel. TSI1 gives $z=0$ in $k'$, contradicting
$z_{g_0}=\bar\Psi(g_0)^{-1}\ne0$. Thus $(M_g)_{20}=0$ for all $g$.
% END TSI2

% BEGIN TSI3
\paragraph{TSI3. Two nonzero weight signs force the original type-I condition.}
A marked characteristic-zero point is represented, after torus
rescaling, by an integral point as in VSCM8. An entry of weight $j$
is multiplied by $t^j$ for one $t\ne0$; vanishing and both
triangularity tests survive this rescaling. The marking at $c_*,g_1$
also survives, since each prescribed difference has one weight.
The diagonal marking difference $\zeta$ has weight zero and residual
value $\bar\Psi(g_1)-\bar\Psi(g_1)^{-1}$. This is nonzero
by the marked-unit condition. Apply the marked shape classification SHP2
to the given point over $\overline{\mathbb Q}_\ell$. The two
triangular alternatives contradict the two nonzero signs (VSCM9).
In the remaining nontype-I alternative the $(2,1)$ entry vanishes
on the sign kernel, and the block on coordinates $0,2$ is irreducible
there. Transfer the entry vanishing to the integral point and apply
TSI2. Transfer the resulting $(2,0)$ vanishing back. The line
spanned by the first coordinate of that block is now invariant,
nonzero and proper, a contradiction. This proves the required type-I
test on the native geometric points.
% END TSI3

% BEGIN HVJ1
\paragraph{HVJ1. The native count and length on the same Case V frame.}
At a host satisfying the guard of CVI1, construct $S,A,F$ and its
arithmetic frame witness as in CVI2. Keep the same $F$ in every
step below. Use exactly the counting hypotheses: Lemmas M and M-prime,
the local finiteness, duality and Euler-characteristic statements,
global finiteness and Euler characteristic, normalized Poitou--Tate,
and the source host inputs for vanishing, weight $-2$, rigidity,
separation and the strengthened marked theorem.

The frame has the residual vanishing and five Euler defects of CVI3--4.
The strengthened marked theorem supplies $g_1$ and finiteness on the
native universal pair. NMD1--3 construct this pair with its integral
origin, NMD4--7 bound its tangent and relation multiplicities, and VSCM4
proves the nonpositive integral tangent test from the host vanishing
hypothesis. VSCM5--11 identify its graded entry algebra, geometric
points, sign conditions, isomorphism classes and derivations. In
particular a point vanishing on both signs is the diagonal origin by
the class-one condition for this $F$. At a point with both signs
TSI3 gives type I, using the host weight $-2$ input for the very
$X$-data supplied by CVI1. Thus the strengthened theorem gives the
prescribed automorphy and irreducibility record. The host rigidity and
number hypotheses apply to that same point. VSCM10 carries rigidity
to the Euler-derivation test. The separating function of degree zero
is exactly the trace difference at the single separating element $\phi$.

These supply all hypotheses of CVI5 on this model. If there are no
weight $-1$ variables, VSCM1--3 inject each finite Selmer group into
the zero module and one takes $V=0$ and an empty list. Otherwise VSCM11
and NMD7 give the homogeneous algebra and five padded equation counts
of CVI5. The original valuation and length lemmas give a finite
irredundant list of two-sided points and an integer $V$ bounding the
length of the solution module in weight $-1$ modulo each nonzero $b$.
VSCM1--3 inject the native finite Selmer module into those solutions
on this very ring. VSCM8 realizes the points integrally after scaling,
which leaves their degree-zero traces unchanged. VSCM6--7 identify the
orbit and trace comparisons used to keep the list irredundant.
Hence on $F$ the same list satisfies
\[
  \operatorname{length}_O(\mathrm{Sha}_b)\le V,\qquad
  V/e_O\le2\sum_i v_\ell\!\left(
    \operatorname{tr} r_i(\phi)-\operatorname{tr}r_z(\phi)\right),
  \qquad e_O=v_O(\ell).
\]
Here $\mathrm{Sha}_b$ denotes the weight $-1$ finite Selmer module
with coefficients $O/(b)$ used by the source. VLT1 identifies the
original strict group, by inflation, with the strict divisible-coefficient
group on $\Gamma$. By VLT2 its $b$-torsion lies in the image of
$\mathrm{Sha}_b$; away from $T_b$ the finite lift need only vanish on
$N_v$. Thus each such torsion submodule has length at most $V$.
Their exhaustion in VLT3 gives the original $\mathrm{Vlen}\le V$. The coefficient ring in its normalization is
still $O$. This proves the whole host, frame and length comparison
on the stated arithmetic inputs.
% END HVJ1

% BEGIN HVJ2
\paragraph{HVJ2. The original guarded count with its uniform exponent.}
Choose the Gaussian field first. The original host number hypothesis
supplies $C$ before the character and host instances in CVI6.
For each such instance CVI1 supplies both host predicates under
$h.\mathrm{CaseV}(\ell,e)$ and $R_1\le 2^{c_2(\psi)}$, including
equality. Choose $g_1$ as above and the source separating $\phi$ before
the list. Apply the number bound at this $\phi$ to the list in HVJ1:
its representations have the prescribed strengthened automorphy
property and are pairwise inequivalent on the sign kernel. With
$X=\ell N e\ge2$, the source normalization gives
\[
 \mathrm{Vn}\log\ell
 \le (V/e_O)\log\ell
 \le 2\sum_i v_\ell\bigl(\operatorname{tr}r_i(\phi)
                      -\operatorname{tr}r_z(\phi)\bigr)\log\ell
 \le 2X^C\le X^{C+1}.
\]
Together with finiteness from HVJ1 this proves the literal
$\mathrm{CountV}$ in the implication
$\mathrm{countV\_of\_finiteness\_duality\_euler\_host}$.
The exponent works at all allowed hosts and exponents in that implication.
The strengthened marked theorem and the arithmetic host suppliers are
inputs to it.
% END HVJ2
